\Ch{Declarations}{Decl}
\Sec{Preamble}{Decl.Pre}
\p Declarations generally specify how names are to be interpreted. Declarations have the form
\begin{grammar}
  \define{declaration-seq}\br
  declaration\br
  declaration-seq declaration\br

  \define{declaration}\br
  name-declaration\br
  special-declaration\br

  \define{name-declaration}\br
  block-declaration\br
  function-definition\br
  template-declaration\br
  namespace-definition\br
  empty-declaration\br
  attribute-declaration\br
  cbuffer-declaration\br

  \define{special-declaration}\br
  export-declaration-group\br
  cbuffer-member-declaration\br

  \define{block-declaration}\br
  simple-declaration\br
  namespace-alias-definition\br
  using-declaration\br
  using-directive\br
  static\_assert-declaration\br
  alias-declaration\br
  opaque-enum-declaration\br

  \define{empty-declaration}\br
  \terminal{;}

\end{grammar}

\Sec{Specifiers}{Decl.Spec}
\Sub{General}{Decl.Spec.General}
\p The specifiers that can be used in a declaration are
\begin{grammar}
  \define{decl-specifier}\br
  \textit{storage-class-specifier}\br
  \textit{type-specifier}\br
  \textit{function-specifier}\br
  \texttt{typedef}\br

  \define{decl-specifier-seq}\br
  decl-specifier\br
  decl-specifier decl-specifier-seq\br
\end{grammar}

\Sub{Function specifiers}{Decl.Spec.Fct}

\p A \textit{function-specifier} can be used only in a function declaration.

\begin{grammar}
  \define{function-specifier}\br
  \texttt{export}\br
\end{grammar}

\p The \texttt{export} specifier denotes that the function has program linkage (\ref{Basic.Linkage.Program}).

\p The \texttt{export} specifier cannot be used on functions directly or indirectly within an unnamed namespace.

\p Functions with program linkage can also be specified in \textit{export-declaration-group} (\ref{Decl.Export}).

\p If a function is declared with an \texttt{export} specifier then all redeclarations of the same function must also use the \texttt{export} specifier or be part of \textit{export-declaration-group} (\ref{Decl.Export}).

\Sec{Declarators}{Decl.Decl}

\Sec{Declarators}{Decl.Decl}

\begin{grammar}
  \define{init-declarator-list}\br
  init-declarator\br
  init-declarator-list \terminal{,} init-declarator\br

  \define{init-declarator}\br
  declarator \opt{initializer}\br

  \define{declarator}\br
  declarator-id \opt{attribute-specifier-seq}\br
  declarator parameters-and-qualifiers\br
  declarator \terminal{\lbrack} \opt{constant-expression} \terminal{\rbrack} \opt{attribute-specifier-seq}\br

  \define{parameters-and-qualifiers}\br
  \terminal{(} parameter-declaration-clause \terminal{)} \opt{cv-qualifier-seq} \opt{attribute-specifier-seq}\br

  \define{cv-qualifier-seq}\br
  cv-qualifier \opt{cv-qualifier-seq}\br

  \define{cv-qualifier}\br
  \terminal{const}\br

  \define{declarator-id}\br
  id-expression

\end{grammar}

\Sub{Arrays}{Decl.Array}

\p In a declaration \texttt{T D} where \texttt{D} has the form:

\begin{innergrammar}
  \terminal{D1} \terminal{\lbrack} \opt{constant-expression} \terminal{\rbrack} \opt{attribute-specifier-seq}
\end{innergrammar}

\p The type of the entity declared in a declaration \texttt{T D1} is
\texttt{T$^\prime$}, the derived type \texttt{T}. The type of the entity
declared in the declarator \texttt{D} is \texttt{T$^\prime$[N]}. The
\textit{constant-expression} shall be a converted constant expression of
unsigned integer type. Its value \texttt{N} specifies the number of elements in
the array, called the \textit{array bound}.

\p If present, \texttt{N} must be greater than zero. If not present, the
declaration must be a global declaration of an object of a resource type
(\ref{Resources}); otherwise the program is ill-formed.

\Sec{Initializers}{Decl.Init}

\p The process of initialization described in this section applies to all
initializers regardless of the context.

\begin{grammar}
  \define{initializer}\br
  brace-or-equal-initializer\br
  \terminal{(} expression-list \terminal{)}\br

  \define{brace-or-equal-initializer}\br
  \terminal{=} initializer-clause\br
  braced-init-list\br

  \define{initializer-clause}\br
  assignment-expression\br
  braced-init-list\br

  \define{braced-init-list}\br
  \terminal{\{} initializer-list \opt{\terminal{,}} \terminal{\}}\br
  \terminal{\{} \terminal{\}}\br

  \define{initializer-list}\br
  initializer-clause\br
  initializer-list \terminal{,} initializer-clause\br
\end{grammar}

\Sub{Aggregate Initialization}{Decl.Init.Agg}

\p An \textit{aggregate} is a vector, matrix, array, or class.

\p The subobjects of an aggregate have a defined order. For vectors and arrays
the order is increasing subscript order. For matrices it is increasing subscript
order with the subscript nesting such that in the notation
\texttt{Mat[M][N]}, the ordering is \(Mat[0][0]...Mat[0][N]...
Mat[M][0]...Mat[M][N]\). For classes the order is base class, followed by member
subobjects in declaration order.

\p A \textit{flattened ordering} of subobjects can be produced by performing a
depth-first traversal of the subobjects of an object following the defined
subobject ordering.

\p Each \textit{braced initializer list} is comprised of zero or more
\textit{initializer-clause} expressions, which is either another braced
initializer list or an expression which generates a value that either is or can
be implicitly converted to an rvalue. Each assignment-expression is an object,
which may be a scalar or aggregate type. A \textit{flattened initializer
sequence} is a sequence of expressions constructed by a depth-first traversal
over each assignment-expression in an initializer-list and performing a
depth-first traversal accessing each subobject of the assignment-expression.

\p If the target object is an array of unknown size, the object is assumed to
have \(m\) possible elements during parsing, where \(m>0\).

\p An initializer-list is a valid initializer if for each element
\(E_{n \bmod m}\) in the target object's flattened ordering there is a
corresponding expression \(E_n\) in the flattened initializer sequence, which
can be implicitly converted to the element's type. For arrays of unknown size,
the total number of expressions in the flattened initializer sequence must be a
multiple of the array's base element type.

\p An initializer-list is invalid if the flattened initializer sequence contains
more or fewer elements than the target object's flattened ordering, or if any
initializer \(I_n\) cannot be implicitly converted to the corresponding element
\(E_n\)'s type.
